\chapter{Homodyne detection in quantum optics}\label{app:homodyne}
As explained in Section~\ref{sec:homodyne}, the homodyne detection creates a continuous signal containing information on the quadrature of the physical system of interest. We use the notations of Figure~\ref{fig:homodyne}, i.e., the system field is labelled by the system annihilation operator $\hat{a}$ (and is now an input of the homodyne detection process). We also use another input signal which is a coherent field labelled by its annihilation operator $\hat{b}$. Moreover, we use a 50:50 beam-splitter to form two output fields labelled $\hat{c}$ and $\hat{d}$ which are then measured by two photo-detectors, as depicted in the aforementioned figure. The quadrature of the system is defined by
\begin{equation}
    \hat{a}_\theta=\frac{\hat{a}^\dagger e^{i\theta}+\hat{a}e^{-i\theta}}{\sqrt{2}}.
\end{equation}

Quantum optics~\cite{Gerry2023} assures that, after the beam-splitter, the output fields are 
\begin{equation}
    \hat{c}=\frac{\hat{a}+i\hat{b}}{\sqrt{2}},\qquad
    \hat{d}=\frac{\hat{b}+i\hat{a}}{\sqrt{2}}.
\end{equation}
These relations hold when the beam-splitter is 50:50. Hence, this exact process is called ``balanced'' homodyne detection. We explained in Section~\ref{sec:photodet} that the photo-detector signal is proportional to the mean value of the occupation number. Thus, $I_c \propto \expval{\hat{c}^\dagger \hat{c}}$ and $I_d \propto \expval{\hat{d}^\dagger \hat{d}}$. Therefore, the difference of photo-currents is
\begin{equation}\label{eq:D2}
    I_d-I_c \propto i\expval{\hat{a}^\dagger\otimes\hat{b}-\hat{a}\otimes\hat{b}^\dagger}.
\end{equation}
As explained, the field input labelled by $\hat{b}$ is a coherent field. We denote this field as $\ket{\beta}$ so that $\hat{b}\ket{\beta}=\beta\ket{\beta}$, where $\beta=\abs{\beta}e^{i\phi}\in\mathbb{C}$. If we denote the input state of the system as $\ket{\psi}$, then the total input state is $\ket{\psi_\text{in}}=\ket{\psi} \otimes \ket{\beta}$. We can therefore explicit the average in Eq.~(\ref{eq:D2}) as
\begin{equation}
    I_d-I_c \propto i\abs{\beta}\expval{(\hat{a}^\dagger e^{i\phi}-\hat{a}e^{-i\phi})}{\psi} = \abs{\beta}\expval{(\hat{a}^\dagger e^{i\theta}+\hat{a}e^{-i\theta})}{\psi},
\end{equation}
where we introduced $\theta=\phi-\pi$. Therefore, we have briefly shown that homodyne detection signal is proportional to the quadrature of the system $\hat{a}_\theta$.